% !TeX program = lualatex
% Anschreiben (Deutsch), Stelle: OT-C IT-Projektmitarbeiter (OPTITOOL GmbH)
% Design passend zu antonrechenauer_onepage_cv_admin.tex
%
% kompilieren:
%   lualatex Anschreiben_optitool.tex

% ---------------------------------------------------------------- Parameter (Standardwerte)
% leer lassen = Zeile bzw. Satz entfällt
\providecommand*{\Firma}{OPTITOOL GmbH}
\providecommand*{\Abteilung}{Personalabteilung}
\providecommand*{\Ansprechpartner}{}                       % z.\,B. Frau Erika Mustermann
\providecommand*{\Strasse}{Im Gewerbepark B 38}                               % Straße und Hausnummer
\providecommand*{\Ort}{D-93059 Regensburg}
\providecommand*{\Anrede}{Hallo liebes OPTITOOL-Team,}
\providecommand*{\Stelle}{OT-C IT-Projektmitarbeiter (m/w/d)}
\providecommand*{\Referenz}{}                              % Kennziffer der Ausschreibung
\providecommand*{\Quelle}{}                                % z.\,B. auf StepStone / auf Ihrer Karriereseite
\providecommand*{\Eintritt}{}                              % z.\,B. zum 1.~Januar 2027
\providecommand*{\Gehalt}{}                                % z.\,B. 65.000~€ brutto jährlich
\providecommand*{\Datum}{\today}

\documentclass[11pt,a4paper]{article}

\usepackage[ngerman]{babel}
\usepackage{fontspec}
\setmainfont{Lato}[
  Path           = /usr/share/fonts/truetype/lato/,
  Extension      = .ttf,
  UprightFont    = Lato-Regular,
  BoldFont       = Lato-Bold,
  ItalicFont     = Lato-Italic,
  BoldItalicFont = Lato-BoldItalic,
  FontFace       = {sb}{n}{Lato-Semibold},
  FontFace       = {l}{n}{Lato-Light},
  FontFace       = {k}{n}{Lato-Black},
  Numbers        = Lining]
\newcommand*{\semibold}{\fontseries{sb}\selectfont}
\newcommand*{\light}{\fontseries{l}\selectfont}
\newcommand*{\heavy}{\fontseries{k}\selectfont}

\newlength{\bandheight}\setlength{\bandheight}{2.75cm}

\usepackage[a4paper,top=3.6cm,bottom=1.2cm,left=2.5cm,right=2cm]{geometry}
\usepackage{xcolor}
\usepackage{etoolbox}
\usepackage{tikz}
\usepackage{eso-pic}
\usepackage[hidelinks]{hyperref}

\definecolor{band}{HTML}{16222E}
\definecolor{accent}{HTML}{1E8C8A}
\definecolor{ink}{HTML}{1C2530}
\definecolor{muted}{HTML}{5E6A75}
\definecolor{bandtext}{HTML}{B9C6D2}

\hypersetup{
  pdftitle    = {Anton Rechenauer -- Anschreiben \Stelle},
  pdfauthor   = {Anton Rechenauer},
  pdfsubject  = {Bewerbung als \Stelle\ bei \Firma},
  pdfkeywords = {IT-Projektmitarbeiter, Integrationsprojekte, Transportlogistik,
                 Disposition, Flottenmanagement, Customizing, Installation,
                 Konfiguration, Systemumgebungen, Rollout, MS SQL Server, MSSQL,
                 PostgreSQL, T-SQL, Datenstrukturen, Schnittstellen, ETL,
                 Dokumentation, Projektmanagement, Linux, Windows Server, Python},
  pdflang     = {de-DE}}

\pagestyle{empty}
\setlength{\parindent}{0pt}
\setlength{\parskip}{0.55em}
\linespread{1.07}
\color{ink}

% ---------------------------------------------------------------- Kopfband
\AddToShipoutPictureBG*{%
  \begin{tikzpicture}[remember picture,overlay]
    \fill[band] (current page.north west) rectangle
                ([yshift=-\bandheight]current page.north east);
    \fill[accent] ([yshift=-\bandheight]current page.north west) rectangle
                  ([yshift=-\bandheight-2.2pt]current page.north east);
    \node[anchor=north west,inner sep=0pt,text width=17cm,align=left]
      at ([xshift=2.5cm,yshift=-0.55cm]current page.north west) {%
        {\fontsize{24}{26}\selectfont\color{white}\light Anton \heavy Rechenauer}\par
        \vspace{4pt}
        {\fontsize{11.5}{13}\selectfont\color{accent!55!white}\semibold
          IT-Projekte \& Anwendungssupport\,\,{\color{bandtext}\light|}\,\,
          {\light\color{bandtext} Integration · Datenbanken · Customizing}}\par
        \vspace{6pt}
        {\small\color{bandtext}
          Silberne-Fisch-Gasse 16, 93047 Regensburg
          \enspace{\color{accent}/}\enspace +49~1525~3612827
          \enspace{\color{accent}/}\enspace \href{mailto:antonrechenauer@gmail.com}{antonrechenauer@gmail.com}}};
  \end{tikzpicture}}

% Zeile nur ausgeben, wenn der Parameter nicht leer ist
\newcommand*{\optline}[1]{\ifdefempty{#1}{}{#1\\}}

\begin{document}

% ---------------------------------------------------------------- Empfänger + Datum
{\footnotesize\color{muted}Anton Rechenauer · Silberne-Fisch-Gasse 16 · 93047 Regensburg}\par
\vspace{2pt}
\begin{minipage}[t]{0.6\textwidth}
  \raggedright
  {\semibold\Firma}\\
  \optline{\Abteilung}%
  \optline{\Ansprechpartner}%
  \optline{\Strasse}%
  \Ort
\end{minipage}%
\hfill
\begin{minipage}[t]{0.35\textwidth}
  \raggedleft
  \vspace{0pt}\par
  Regensburg, \Datum
\end{minipage}

\vspace{1.0cm}

% ---------------------------------------------------------------- Betreff
{\large\color{band}\heavy Bewerbung als \Stelle}\par
\vspace{-0.4em}
{\small\color{muted}%
  \ifdefempty{\Referenz}{Integrationsprojekte · Installation \& Customizing · MS SQL Server · PostgreSQL}{Kennziffer: \Referenz}}\par

\vspace{0.6em}

% ---------------------------------------------------------------- Text
\Anrede

Sie suchen jemanden, der Integrationsprojekte technisch von der Planung bis zum Go-live
begleitet -- Systemumgebungen aufsetzt, installiert, konfiguriert, Anpassungen ausrollt und
dem Projektmanager dabei den Rücken freihält. Genau das habe ich in den letzten Jahren in
ganz unterschiedlichen Umgebungen gemacht.%
\ifdefempty{\Quelle}{}{ Auf Ihre Ausschreibung bin ich \Quelle\ aufmerksam geworden.}

\textbf{Systeme aufsetzen und am Laufen halten.} Zuletzt war ich bei LyondellBasell in der
Raffinerie Münchsmünster für Administration und IT-Projekte im Prozessleitnetz zuständig:
Windows Server, Active Directory, vSphere und Hyper-V, Firewalls, SQL-Server-Administration
sowie Installation und Konfiguration neuer Hard- und Softwarelösungen -- von der Anforderung
bis zur Umsetzung. Linux-Server (Debian, Apache, Tomcat) betreue ich seit vielen Jahren per SSH.

\textbf{Datenbanken und Datenstrukturen.} Mit MS SQL Server arbeite ich seit über 20~Jahren
als Entwickler (.NET, T-SQL, Access VBA, Analysis Services) und Administrator, PostgreSQL habe
ich als Lead Scheduler für einen Anlagenstillstand administriert und Schnittstellen sowie
Python-Automatisierung dazu gebaut. Ich habe ETL-Werkzeuge entwickelt, die rund 100
verschiedene Quelldateien prüfen, bereinigen und in eine SQL-Datenbank importieren, und mich
in so manche fremde Datenstruktur hineingedacht. Customizing heißt für mich: erst verstehen,
wie der Kunde arbeitet, dann sauber konfigurieren.

\textbf{Projekte, Kommunikation, Dokumentation.} Als Lead Scheduler habe ich ein
großes Projekt geplant, das in Zeit und Budget abgeschlossen wurde, als Projektleiter in der
Automobilzulieferung Lieferanten und OEM koordiniert, Workshops vorbereitet und
Architektur-Reviews organisiert. Lasten- und Pflichtenhefte schreibe ich gern, und als
ehemaliger Gymnasiallehrer und Dozent erkläre ich Technik so, dass sie auch beim Kunden
ankommt -- auf Deutsch wie auf Englisch.

Das Thema Transportlogistik ist mir nicht fremd: Ich habe als Produktmanager in
Produktionsplanung und Logistik angefangen und später am Navigationsmodul einer
Infotainment-Einheit mitgearbeitet. Als Diplom-Mathematiker reizt mich Optimierung ohnehin --
eine Disposition-Software mit integriertem Flottenmanagement ist da ein sehr spannendes
Produkt. Und dass Sie in Regensburg sitzen, passt perfekt: Ich wohne hier.%
\ifdefempty{\Eintritt}{}{ Für Sie stehe ich \Eintritt\ zur Verfügung.}%
\ifdefempty{\Gehalt}{}{ Meine Gehaltsvorstellung liegt bei \Gehalt.}
Ich freue mich darauf, Sie kennenzulernen!

Mit freundlichen Grüßen

\vspace{1.0cm}
Anton Rechenauer

\vfill
{\small\color{muted}\semibold Anlagen:\enspace\light Lebenslauf, Zeugnisse}

\end{document}
